% ============================================================
%  GUÍA 1: EMPEZANDO A TRABAJAR CON ÁNGULOS
%  Curso 4to Año — Matemática
%  Autor: Ing. Gabriel Astudillo
%  Clases cubiertas: 1 a 3
% ============================================================

\documentclass[12pt, a4paper]{article}

% ── Paquetes ─────────────────────────────────────────────────

\usepackage{mathwizards-guia}
\mwguia{Guía 1 — Ángulos}{Ing. Gabriel Astudillo $\cdot$ 4to Año}

\begin{document}
% ============================================================

% ── Portada / Título ─────────────────────────────────────────
\mwportada{Guía de Estudio 1}{Empezando a trabajar con ángulos}{Concepto, clasificación, sistemas de medición y ángulos en el plano cartesiano}

\vspace{4pt}
\begin{cuadroinfo}
  \small
  \textbf{Presentación:} ¡Hola! Los ángulos están en todas partes: en las manecillas del reloj, en las esquinas de tu cuaderno, en las rampas y hasta en el lanzamiento de una pelota. En esta guía vamos a convertirnos en \emph{cazadores de ángulos}: aprenderemos a reconocerlos, clasificarlos, medirlos en grados y en radianes, y a ubicarlos en el plano cartesiano. Cada sección tiene una explicación sencilla, ejemplos resueltos paso a paso y ejercicios que van de lo más fácil a verdaderos retos. ¡Tú puedes!
  \\[4pt]
  \textbf{Nivel de Dificultad:}\quad
  \facil{} Básico / Directo \quad
  \medio{} Intermedio \quad
  \dificil{} Desafiante
\end{cuadroinfo}

\vspace{8pt}

% ============================================================
\section{Concepto, clasificación y relaciones entre ángulos}
% ============================================================

\subsection{¿Qué es un ángulo?}

\begin{cuadroteoria}
  \textbf{Ángulo:} es la abertura formada por dos \textbf{semirrectas} (los \textbf{lados}) que comparten un punto común llamado \textbf{vértice}.
  \begin{center}
    \begin{tikzpicture}[scale=1]
      \draw[thick, -Latex] (0,0) -- (3.2,0) node[below right] {$B$};
      \draw[thick, -Latex] (0,0) -- (1.6,2.8) node[above] {$A$};
      \draw[mwprimario, thick] (0.7,0) arc (0:60:0.7);
      \node[mwprimario, font=\small] at (1.05,0.28) {$\alpha$};
      \filldraw (0,0) circle (2pt) node[below left] {$O$};
    \end{tikzpicture}
  \end{center}
  El ángulo de la figura se denota $\angle AOB$ (el vértice va en el medio) o simplemente $\alpha$.

  \textbf{Medida:} los ángulos se miden en \textbf{grados} ($^{\circ}$). Un giro completo mide $360^{\circ}$; un ángulo recto mide $90^{\circ}$.
\end{cuadroteoria}

\newpage %ajuste pag

\begin{cuadroteoria}
  \textbf{Clasificación según su medida:}
  \begin{itemize}[leftmargin=*]
    \item \textbf{Agudo:} mide menos de $90^{\circ}$ (ej. $35^{\circ}$).
    \item \textbf{Recto:} mide exactamente $90^{\circ}$.
    \item \textbf{Obtuso:} mide más de $90^{\circ}$ y menos de $180^{\circ}$ (ej. $120^{\circ}$).
    \item \textbf{Llano:} mide exactamente $180^{\circ}$ (forma una línea recta).
    \item \textbf{Completo:} mide exactamente $360^{\circ}$ (un giro completo).
  \end{itemize}

  \textbf{Clasificación según su posición:}
  \begin{itemize}[leftmargin=*]
    \item \textbf{Adyacentes:} comparten el vértice y un lado (ej. $\angle AOB$ y $\angle BOC$ con lado común $OB$).
    \item \textbf{Opuestos por el vértice:} los lados de uno son prolongación de los lados del otro. \emph{Los ángulos opuestos por el vértice son iguales.}
  \end{itemize}
\end{cuadroteoria}

\begin{cuadroteoria}
  \textbf{Ángulos complementarios y suplementarios:}
  \begin{itemize}[leftmargin=*]
    \item \textbf{Complementarios:} dos ángulos cuyas medidas suman $90^{\circ}$.
          $$\text{complemento de } \alpha = 90^{\circ} - \alpha$$
    \item \textbf{Suplementarios:} dos ángulos cuyas medidas suman $180^{\circ}$.
          $$\text{suplemento de } \alpha = 180^{\circ} - \alpha$$
  \end{itemize}
\end{cuadroteoria}

\begin{cuadroadvertencia}
  \textbf{Errores clásicos:}
  \begin{itemize}[leftmargin=*]
    \item Confundir \emph{complementario} (suma $90^{\circ}$) con \emph{suplementario} (suma $180^{\circ}$). Truco: el \emph{complemento} completa un ángulo \emph{recto}; el \emph{suplemento} completa un ángulo \emph{llano}.
    \item Pensar que dos ángulos adyacentes siempre suman $180^{\circ}$: solo es cierto si están sobre una misma recta.
    \item Olvidar que los ángulos opuestos por el vértice son \emph{iguales} y los adyacentes sobre una recta son \emph{suplementarios}.
  \end{itemize}
\end{cuadroadvertencia}

\newpage %ajuste pag

\subsection{Ejercicios resueltos}

\textbf{Ejemplo R1.} Clasifica cada ángulo según su medida: $a)\ 35^{\circ} \quad b)\ 90^{\circ} \quad c)\ 145^{\circ} \quad d)\ 180^{\circ}$.

\begin{cuadrosolucion}
  \textbf{Solución:}
  \begin{itemize}
    \item $a)\ 35^{\circ}$: agudo (menos de $90^{\circ}$).
    \item $b)\ 90^{\circ}$: recto.
    \item $c)\ 145^{\circ}$: obtuso (entre $90^{\circ}$ y $180^{\circ}$).
    \item $d)\ 180^{\circ}$: llano.
  \end{itemize}
\end{cuadrosolucion}

\textbf{Ejemplo R2.} Calcula el complemento y el suplemento de $42^{\circ}$.

\begin{cuadrosolucion}
  \textbf{Solución:}
  \begin{itemize}
    \item Complemento: $90^{\circ} - 42^{\circ} = 48^{\circ}$.
    \item Suplemento: $180^{\circ} - 42^{\circ} = 138^{\circ}$.
  \end{itemize}
  \emph{Comprobación:} $42^{\circ} + 48^{\circ} = 90^{\circ}$ y $42^{\circ} + 138^{\circ} = 180^{\circ}$. \checkmark
\end{cuadrosolucion}

\textbf{Ejemplo R3.} Dos rectas se cortan formando los ángulos opuestos por el vértice $\alpha = (3x - 10)^{\circ}$ y $\beta = (2x + 15)^{\circ}$. Halla el valor de $x$ y la medida de cada ángulo.

\begin{cuadrosolucion}
  \textbf{Solución:} como son opuestos por el vértice, son iguales:
  $$3x - 10 = 2x + 15 \quad \Rightarrow \quad 3x - 2x = 15 + 10 \quad \Rightarrow \quad x = 25.$$
  Sustituimos: $\alpha = 3(25) - 10 = 75 - 10 = 65^{\circ}$ y $\beta = 2(25) + 15 = 65^{\circ}$. \checkmark
\end{cuadrosolucion}

\subsection{Ejercicios propuestos}

\begin{enumerate}[leftmargin=*, resume]
  \item \facil{} Clasifica según su medida: $a)\ 15^{\circ} \qquad b)\ 90^{\circ} \qquad c)\ 120^{\circ} \qquad d)\ 180^{\circ} \qquad e)\ 360^{\circ}$

  \item \facil{} Clasifica según su posición: dos ángulos que comparten el vértice y un lado se llaman \_\_\_\_\_\_\_\_; dos ángulos cuyos lados son prolongación del otro se llaman \_\_\_\_\_\_\_\_.

  \item \facil{} Calcula el complemento de: $a)\ 30^{\circ} \qquad b)\ 58^{\circ} \qquad c)\ 12^{\circ}$

  \item \facil{} Calcula el suplemento de: $a)\ 70^{\circ} \qquad b)\ 95^{\circ} \qquad c)\ 15^{\circ}$

  \item \facil{} ¿Qué ángulo forman las manecillas del reloj a las 3:00? ¿Y a las 6:00?

  \item \medio{} Dos ángulos son complementarios y uno mide el doble del otro. ¿Cuánto mide cada uno?

  \item \medio{} El suplemento de un ángulo es 4 veces su medida. Halla el ángulo.

  \item \medio{} Dos ángulos opuestos por el vértice miden $(2x + 10)^{\circ}$ y $(3x - 20)^{\circ}$. Halla $x$ y la medida de cada ángulo.

  \item \medio{} Dos ángulos adyacentes sobre una misma recta miden $(x + 15)^{\circ}$ y $(2x - 30)^{\circ}$. Halla $x$ y cada ángulo.

  \item \dificil{} Una transversal corta a dos rectas paralelas. Uno de los ángulos correspondientes mide $55^{\circ}$. ¿Cuánto miden sus alternos internos y los adyacentes a él sobre la misma recta?

  \item \dificil{} Un carpintero necesita cortar una tabla con dos cortes que formen ángulos complementarios, de modo que uno mida $10^{\circ}$ más que el otro. ¿Cuáles son las medidas de los dos cortes?
\end{enumerate}

% ============================================================
\section{Sistemas de medición angular}
% ============================================================

\subsection{El sistema sexagesimal}

\begin{cuadroteoria}
  En el \textbf{sistema sexagesimal} el grado se divide en 60 partes iguales:
  $$1^{\circ} = 60' \quad (\text{minutos}) \qquad 1' = 60'' \quad (\text{segundos})$$
  Por ejemplo, $35^{\circ} 24' 18''$ se lee: 35 grados, 24 minutos y 18 segundos.
\end{cuadroteoria}

\begin{cuadropasos}
  \textbf{Convertir de grados a minutos (o de minutos a segundos):} se multiplica por 60.
  \textbf{Convertir de minutos a grados (o de segundos a minutos):} se divide entre 60.
  \textbf{Pasar de una forma compleja a una simple:} $35^{\circ} 30' = 35^{\circ} + \dfrac{30}{60}^{\circ} = 35{,}5^{\circ}$.
\end{cuadropasos}

\subsection{El radián}

\begin{cuadroteoria}
  En el \textbf{sistema circular} los ángulos se miden en \textbf{radianes} (rad). Un radián es el ángulo central que abarca un arco de longitud igual al radio de la circunferencia.
  \begin{center}
    \begin{tikzpicture}[scale=0.9]
      \draw[thick, mwgris] (0,0) circle (2);
      \draw[thick, -Latex] (0,0) -- (2,0) node[right] {$r$};
      \draw[thick, -Latex, mwprimario] (0,0) -- (1,1.732);
      \draw[mwprimario, thick] (0.5,0) arc (0:60:0.5);
      \node[mwprimario] at (0.85,0.22) {$1 \text{ rad}$};
      \filldraw (0,0) circle (2pt) node[below left] {$O$};
    \end{tikzpicture}
  \end{center}
  \textbf{Equivalencias fundamentales:}
  $$180^{\circ} = \pi \text{ rad} \qquad 360^{\circ} = 2\pi \text{ rad} \qquad 90^{\circ} = \frac{\pi}{2} \text{ rad}$$
\end{cuadroteoria}

\begin{cuadropasos}
  \textbf{Convertir de grados a radianes:} se multiplica por $\dfrac{\pi}{180^{\circ}}$.
  $$45^{\circ} = 45^{\circ} \cdot \frac{\pi}{180^{\circ}} = \frac{45\pi}{180} = \frac{\pi}{4} \text{ rad}$$

  \textbf{Convertir de radianes a grados:} se multiplica por $\dfrac{180^{\circ}}{\pi}$.
  $$\frac{2\pi}{3} \text{ rad} = \frac{2\pi}{3} \cdot \frac{180^{\circ}}{\pi} = 120^{\circ}$$
\end{cuadropasos}

\begin{cuadroadvertencia}
  \textbf{Errores clásicos:}
  \begin{itemize}[leftmargin=*]
    \item Confundir $\pi$ radianes ($180^{\circ}$) con $2\pi$ radianes ($360^{\circ}$).
    \item Olvidar simplificar la fracción al convertir (ej. $\dfrac{45\pi}{180}$ se simplifica a $\dfrac{\pi}{4}$).
    \item Sumar grados, minutos y segundos sin agrupar: en $35^{\circ} 24' + 18^{\circ} 45'$, los minutos suman $69' = 1^{\circ} 9'$, así que el resultado es $53^{\circ} + 1^{\circ} 9' = 54^{\circ} 9'$.
  \end{itemize}
\end{cuadroadvertencia}

\subsection{Ejercicios resueltos}

\textbf{Ejemplo R1.} Expresa en minutos: $a)\ 4^{\circ} \qquad b)\ 2^{\circ} 15'$.

\begin{cuadrosolucion}
  \textbf{Solución:}
  \begin{itemize}
    \item $a)\ 4^{\circ} = 4 \cdot 60' = 240'$.
    \item $b)\ 2^{\circ} 15' = 2 \cdot 60' + 15' = 120' + 15' = 135'$.
  \end{itemize}
\end{cuadrosolucion}

\textbf{Ejemplo R2.} Suma: $35^{\circ} 24' + 18^{\circ} 45'$.

\begin{cuadrosolucion}
  \textbf{Solución:} sumamos por separado:
  $$35^{\circ} 24' + 18^{\circ} 45' = (35 + 18)^{\circ} + (24 + 45)' = 53^{\circ} 69'.$$
  Como $69' = 1^{\circ} 9'$, agrupamos: $53^{\circ} 69' = 54^{\circ} 9'$.
\end{cuadrosolucion}

\textbf{Ejemplo R3.} Convierte a radianes: $a)\ 60^{\circ} \qquad b)\ 150^{\circ}$.

\begin{cuadrosolucion}
  \textbf{Solución:}
  \begin{itemize}
    \item $a)\ 60^{\circ} = 60^{\circ} \cdot \dfrac{\pi}{180^{\circ}} = \dfrac{60\pi}{180} = \dfrac{\pi}{3}$ rad.
    \item $b)\ 150^{\circ} = 150^{\circ} \cdot \dfrac{\pi}{180^{\circ}} = \dfrac{150\pi}{180} = \dfrac{5\pi}{6}$ rad.
  \end{itemize}
\end{cuadrosolucion}

\textbf{Ejemplo R4.} Convierte a grados: $a)\ \dfrac{\pi}{4} \text{ rad} \qquad b)\ \dfrac{5\pi}{3} \text{ rad}$.

\begin{cuadrosolucion}
  \textbf{Solución:}
  \begin{itemize}
    \item $a)\ \dfrac{\pi}{4} \cdot \dfrac{180^{\circ}}{\pi} = \dfrac{180^{\circ}}{4} = 45^{\circ}$.
    \item $b)\ \dfrac{5\pi}{3} \cdot \dfrac{180^{\circ}}{\pi} = \dfrac{5 \cdot 180^{\circ}}{3} = 300^{\circ}$.
  \end{itemize}
\end{cuadrosolucion}

\subsection{Ejercicios propuestos}

\begin{enumerate}[leftmargin=*, resume]
  \item \facil{} Expresa en minutos: $a)\ 3^{\circ} \qquad b)\ 8^{\circ} \qquad c)\ 1^{\circ} 30'$

  \item \facil{} Expresa en grados: $a)\ 120' \qquad b)\ 240' \qquad c)\ 45'$

  \item \facil{} Convierte a radianes: $a)\ 180^{\circ} \qquad b)\ 90^{\circ} \qquad c)\ 360^{\circ}$

  \item \facil{} Convierte a grados: $a)\ \dfrac{\pi}{2} \qquad b)\ \pi \qquad c)\ \dfrac{3\pi}{2}$

  \item \medio{} Suma: $35^{\circ} 24' + 18^{\circ} 45'$ (verifica el ejemplo R2) y resta: $90^{\circ} - 43^{\circ} 27'$.

  \item \medio{} Convierte a radianes: $a)\ 45^{\circ} \qquad b)\ 210^{\circ}$ y a grados: $c)\ \dfrac{5\pi}{6}$.

  \item \medio{} Expresa $30^{\circ} 15'$ en grados con decimales.

  \item \medio{} Un ángulo mide $\dfrac{\pi}{5}$ rad. ¿Cuántos grados son?

  \item \dificil{} Convierte $240^{\circ}$ a radianes y $300^{\circ}$ a radianes. ¿Qué observas con los denominadores?

  \item \dificil{} La rueda de una bicicleta gira $3$ vueltas completas. ¿Cuántos grados giró? ¿Y cuántos radianes? (Una vuelta son $360^{\circ}$.)

  \item \dificil{} Un ángulo mide $72^{\circ}$. ¿Cuánto mide en radianes? Expresa la fracción simplificada y verifica que $\dfrac{2\pi}{5}$ rad sea equivalente.
\end{enumerate}

\newpage

% ============================================================
\section{Ángulos en el plano cartesiano}
% ============================================================

\subsection{Ángulo en posición normal}

\begin{cuadroteoria}
  Un ángulo está en \textbf{posición normal} cuando su vértice está en el \textbf{origen} del plano cartesiano y su lado inicial coincide con el \textbf{eje $X$ positivo}.
  \begin{center}
    \begin{tikzpicture}[scale=0.85]
      \draw[thin, gray!40] (-3.5,-1.5) grid (3.5,3.5);
      \draw[thick, -Latex] (-3.8,0) -- (3.8,0) node[right] {$x$};
      \draw[thick, -Latex] (0,-1.8) -- (0,3.8) node[above] {$y$};
      \draw[gray] (0,0) node[below left] {$0$};
      \draw[thick, -Latex, mwprimario] (0,0) -- (2.6,2.6);
      \draw[mwprimario, thick] (0.8,0) arc (0:45:0.8);
      \node[mwprimario] at (1.1,0.3) {$\theta$};
      \node at (-2.7,2.9) {\small I};
      \node at (2.7,2.9) {\small II};
      \node at (2.7,-1.1) {\small III};
      \node at (-2.7,-1.1) {\small IV};
    \end{tikzpicture}
  \end{center}
  \begin{itemize}[leftmargin=*]
    \item Si gira en \textbf{sentido antihorario} (contrario a las agujas del reloj), el ángulo es \textbf{positivo}.
    \item Si gira en \textbf{sentido horario}, el ángulo es \textbf{negativo}.
  \end{itemize}
\end{cuadroteoria}

\begin{cuadroteoria}
  \textbf{Ángulos coterminales:} dos ángulos son \textbf{coterminales} cuando comparten el mismo lado terminal. Se obtienen sumando o restando $360^{\circ}$ (o $2\pi$ rad):
  $$\theta \pm 360^{\circ}, \quad \theta \pm 720^{\circ}, \ldots$$
  Por ejemplo, $45^{\circ}$ y $405^{\circ}$ son coterminales, y también lo es $-315^{\circ}$.
\end{cuadroteoria}

\begin{cuadropasos}
  \textbf{Determinar el cuadrante de un ángulo:}
  \begin{enumerate}[leftmargin=*]
    \item Si el ángulo es mayor que $360^{\circ}$ o menor que $0^{\circ}$, busca un ángulo coterminal sumando o restando $360^{\circ}$ las veces necesarias.
    \item Con el ángulo resultante entre $0^{\circ}$ y $360^{\circ}$, ubícalo según el cuadrante:
          \begin{itemize}[leftmargin=*]
            \item $0^{\circ} < \theta < 90^{\circ}$: cuadrante I.
            \item $90^{\circ} < \theta < 180^{\circ}$: cuadrante II.
            \item $180^{\circ} < \theta < 270^{\circ}$: cuadrante III.
            \item $270^{\circ} < \theta < 360^{\circ}$: cuadrante IV.
          \end{itemize}
  \end{enumerate}
\end{cuadropasos}
\newpage
\begin{cuadroadvertencia}
  \textbf{Errores clásicos:}
  \begin{itemize}[leftmargin=*]
    \item Confundir los cuadrantes: el cuadrante I es \emph{arriba a la derecha} y se avanza en sentido antihorario (I, II, III, IV).
    \item Olvidar que un ángulo negativo como $-45^{\circ}$ gira \emph{hacia abajo} (equivale a $315^{\circ}$, cuadrante IV).
    \item Creer que $45^{\circ}$ y $405^{\circ}$ son ángulos distintos: son el \emph{mismo} ángulo (coterminales).
  \end{itemize}
\end{cuadroadvertencia}

\subsection{Ejercicios resueltos}

\textbf{Ejemplo R1.} Determina el cuadrante de: $a)\ 120^{\circ} \qquad b)\ 315^{\circ} \qquad c)\ 730^{\circ}$.

\begin{cuadrosolucion}
  \textbf{Solución:}
  \begin{itemize}
    \item $a)\ 120^{\circ}$: está entre $90^{\circ}$ y $180^{\circ}$, cuadrante II.
    \item $b)\ 315^{\circ}$: está entre $270^{\circ}$ y $360^{\circ}$, cuadrante IV.
    \item $c)\ 730^{\circ} = 730^{\circ} - 360^{\circ} = 370^{\circ} - 360^{\circ} = 10^{\circ}$ (restamos dos veces $360^{\circ}$): cuadrante I.
  \end{itemize}
\end{cuadrosolucion}

\textbf{Ejemplo R2.} Escribe dos ángulos coterminales de $135^{\circ}$, uno positivo y uno negativo.

\begin{cuadrosolucion}
  \textbf{Solución:}
  \begin{itemize}
    \item Positivo: $135^{\circ} + 360^{\circ} = 495^{\circ}$.
    \item Negativo: $135^{\circ} - 360^{\circ} = -225^{\circ}$.
  \end{itemize}
  Los tres ángulos comparten el mismo lado terminal.
\end{cuadrosolucion}

\textbf{Ejemplo R3.} El lado terminal de un ángulo en posición normal pasa por el punto $(1, 1)$. ¿Cuánto mide el ángulo?

\begin{cuadrosolucion}
  \textbf{Solución:} el punto $(1, 1)$ está en la recta $y = x$, que forma exactamente $45^{\circ}$ con el eje $X$ positivo. Además $(1, 1)$ está en el cuadrante I, así que el ángulo mide $45^{\circ}$.
  \begin{center}
    \begin{tikzpicture}[scale=0.8]
      \draw[thin, gray!40] (-1.5,-0.5) grid (3.5,3.5);
      \draw[thick, -Latex] (-1.8,0) -- (3.8,0) node[right] {$x$};
      \draw[thick, -Latex] (0,-0.8) -- (0,3.8) node[above] {$y$};
      \draw[thick, -Latex, mwprimario] (0,0) -- (2.9,2.9);
      \draw[mwprimario, thick] (0.9,0) arc (0:45:0.9);
      \node[mwprimario] at (1.2,0.3) {$45^{\circ}$};
      \filldraw[mwprimario] (1,1) circle (2.5pt) node[above right] {$(1,1)$};
    \end{tikzpicture}
  \end{center}
\end{cuadrosolucion}

\newpage %ajuste para pág

\subsection{Ejercicios propuestos}

\begin{enumerate}[leftmargin=*, resume]
  \item \facil{} Determina el cuadrante de: $a)\ 45^{\circ} \qquad b)\ 120^{\circ} \qquad c)\ 240^{\circ} \qquad d)\ 315^{\circ}$

  \item \facil{} ¿En qué cuadrante está $150^{\circ}$? ¿Y $-30^{\circ}$? (Dibuja si te ayuda.)

  \item \facil{} Halla un ángulo coterminal positivo de $380^{\circ}$.

  \item \medio{} Determina el cuadrante de: $a)\ 730^{\circ} \qquad b)\ -200^{\circ} \qquad c)\ 500^{\circ}$

  \item \medio{} El lado terminal de un ángulo en posición normal pasa por el punto $(-2, 2)$. ¿Cuánto mide el ángulo?

  \item \medio{} Escribe dos ángulos coterminales de $135^{\circ}$, uno positivo y uno negativo.

  \item \medio{} Un ángulo mide $-120^{\circ}$. Exprésalo como un ángulo positivo equivalente entre $0^{\circ}$ y $360^{\circ}$ e indica su cuadrante.

  \item \dificil{} El lado terminal de un ángulo en posición normal pasa por el punto $(\sqrt{3}, 1)$. ¿Cuánto mide el ángulo? (Ayuda: compara con el ángulo de $30^{\circ}$.)

  \item \dificil{} ¿Qué ángulo, en radianes, recorre la manecilla de los minutos de un reloj en 15 minutos?

  \item \dificil{} Un dron gira en círculos y su orientación actual es $-270^{\circ}$. Exprésala como ángulo positivo entre $0^{\circ}$ y $360^{\circ}$ e indica su cuadrante.

  \item \dificil{} \textbf{El reto de los campeones:} ¿Qué ángulo forman las manecillas del reloj a las 7:20? (Pista: la manecilla de las horas no apunta exactamente a las 7: a las 7:20 ha avanzado $\frac{20}{60} \cdot 30^{\circ} = 10^{\circ}$ más allá de las 7.)
\end{enumerate}

\newpage

% ============================================================
\section{Clave de Respuestas}
% ============================================================

\subsection*{Sección 1: Concepto, clasificación y relaciones (Ejercicios 1--11)}

\begin{enumerate}[leftmargin=*, label=\textbf{\arabic*.}]
  \setcounter{enumi}{0}
  \item a) Agudo. \quad b) Recto. \quad c) Obtuso. \quad d) Llano. \quad e) Completo.
  \item Adyacentes; opuestos por el vértice.
  \item a) $60^{\circ}$. \quad b) $32^{\circ}$. \quad c) $78^{\circ}$.
  \item a) $110^{\circ}$. \quad b) $85^{\circ}$. \quad c) $165^{\circ}$.
  \item A las 3:00 forman $90^{\circ}$; a las 6:00 forman $180^{\circ}$.
  \item Si $\alpha + \beta = 90^{\circ}$ y $\beta = 2\alpha$, entonces $3\alpha = 90^{\circ}$: los ángulos son $30^{\circ}$ y $60^{\circ}$.
  \item $\alpha + 4\alpha = 180^{\circ} \Rightarrow \alpha = 36^{\circ}$.
  \item $2x + 10 = 3x - 20 \Rightarrow x = 30$; cada ángulo mide $70^{\circ}$.
  \item $(x+15) + (2x-30) = 180 \Rightarrow x = 65$; los ángulos miden $80^{\circ}$ y $100^{\circ}$.
  \item Los alternos internos miden $55^{\circ}$; los adyacentes miden $125^{\circ}$.
  \item Los ángulos miden $40^{\circ}$ y $50^{\circ}$ (pues $x + (x+10) = 90$).
\end{enumerate}

\subsection*{Sección 2: Sistemas de medición angular (Ejercicios 12--22)}

\begin{enumerate}[leftmargin=*, label=\textbf{\arabic*.}]
  \setcounter{enumi}{11}
  \item a) $180'$. \quad b) $480'$. \quad c) $90'$.
  \item a) $2^{\circ}$. \quad b) $4^{\circ}$. \quad c) $0{,}75^{\circ}$.
  \item a) $\pi$ rad. \quad b) $\dfrac{\pi}{2}$ rad. \quad c) $2\pi$ rad.
  \item a) $90^{\circ}$. \quad b) $180^{\circ}$. \quad c) $270^{\circ}$.
  \item Suma: $54^{\circ} 9'$. Resta: $90^{\circ} - 43^{\circ} 27' = 46^{\circ} 33'$.
  \item a) $\dfrac{\pi}{4}$ rad. \quad b) $\dfrac{7\pi}{6}$ rad. \quad c) $150^{\circ}$.
  \item $30^{\circ} 15' = 30{,}25^{\circ}$.
  \item $\dfrac{\pi}{5} \cdot \dfrac{180^{\circ}}{\pi} = 36^{\circ}$.
  \item $240^{\circ} = \dfrac{4\pi}{3}$ rad; $300^{\circ} = \dfrac{5\pi}{3}$ rad. Ambos tienen denominador 3.
  \item $3$ vueltas $= 1080^{\circ} = 6\pi$ rad.
  \item $72^{\circ} = \dfrac{72\pi}{180} = \dfrac{2\pi}{5}$ rad. \checkmark
\end{enumerate}

\subsection*{Sección 3: Ángulos en el plano cartesiano (Ejercicios 23--33)}

\begin{enumerate}[leftmargin=*, label=\textbf{\arabic*.}]
  \setcounter{enumi}{22}
  \item a) I. \quad b) II. \quad c) III. \quad d) IV.
  \item $150^{\circ}$: cuadrante II. $-30^{\circ}$: equivale a $330^{\circ}$, cuadrante IV.
  \item $380^{\circ} - 360^{\circ} = 20^{\circ}$.
  \item a) $730^{\circ} = 10^{\circ}$: I. \quad b) $-200^{\circ} = 160^{\circ}$: II. \quad c) $500^{\circ} = 140^{\circ}$: II.
  \item El punto $(-2, 2)$ está en la recta $y = -x$ del cuadrante II: el ángulo mide $135^{\circ}$.
  \item $495^{\circ}$ y $-225^{\circ}$ (cualquier ángulo que difiera en $360^{\circ}$ es válido).
  \item $-120^{\circ} + 360^{\circ} = 240^{\circ}$: cuadrante III.
  \item El punto $(\sqrt{3}, 1)$ corresponde a $30^{\circ}$ (cuadrante I).
  \item $15$ min $= \dfrac{1}{4}$ de vuelta: $\dfrac{2\pi}{4} = \dfrac{\pi}{2}$ rad.
  \item $-270^{\circ} + 360^{\circ} = 90^{\circ}$: no está en ningún cuadrante, está sobre el eje $Y$ positivo.
  \item \textbf{Reto:} la manecilla de los minutos está en el 4 ($120^{\circ}$); la de las horas avanza $7 \cdot 30^{\circ} + 10^{\circ} = 220^{\circ}$. El ángulo menor es $220^{\circ} - 120^{\circ} = 100^{\circ}$.
\end{enumerate}

\vspace{8pt}
\begin{cuadroinfo}
  \small
  \textbf{¿Cómo te fue?} Si lograste resolver hasta el reto de los campeones, ¡ya estás listo para ver los ángulos en acción! En las próximas guías los ángulos aparecerán en la física: en los \emph{vectores}, cuando una fuerza empuje en diagonal, y en el \emph{movimiento} de los cuerpos. ¿Nos vemos ahí?
\end{cuadroinfo}

\end{document}
